\section{Problemas de calidad}

Durante la ejecución de las estructuras de la planta pueden presentarse diversos problemas de calidad, sea por deficiencias en la colocación y compactación, por las condiciones ambientales extremas del sector de La Negra, o por interrupciones no previstas del hormigonado. A continuación se identifican tres problemas representativos de las estructuras estudiadas (muro y losa de fondo del CS, y losa de fundación del campamento), indicando en cada caso su causa, la medida preventiva que pudo adoptarse en el proceso constructivo y la metodología de reparación que se utilizaría. Los tres casos seleccionados corresponden a los tres tipos de reparación revisados en el curso: reconstitución del hormigón, reparación de fisuras y detención de filtraciones \cite{ACI546}.

\subsection{Nidos de piedra (coqueras) en el muro del CS}

\paragraph{Descripción y causa}
Los nidos de piedra, o coqueras, corresponden a zonas localizadas donde el árido grueso queda expuesto y sin pasta de cemento que lo envuelva, dejando una matriz porosa y sin continuidad. En el muro del CS este defecto es particularmente probable, dado que se trata de una sección delgada (35 cm) con alta densidad de armadura, condición que dificulta la penetración del vibrador. Sus causas habituales son una altura libre de caída excesiva que induce segregación, una compactación insuficiente o mal distribuida en las zonas inferiores y congestionadas, y un desplazamiento horizontal del hormigón con el vibrador. En una estructura de contención de líquidos este defecto es crítico, ya que constituye una vía directa de filtración y un punto débil frente al ataque químico de las aguas residuales.

\paragraph{Medida preventiva en el proceso constructivo}
La aparición de coqueras se previene limitando la altura libre de caída del hormigón a un máximo de 1,5 m, colocando el hormigón en capas horizontales de 30 a 40 cm y respetando los parámetros de vibración interna especificados (radio de acción, espaciado de inserción y velocidad de retiro). Para el muro del CS, sin embargo, la medida preventiva más efectiva es la ya adoptada en la especificación: el uso de hormigón autocompactante (HAC), que se consolida por gravedad sin vibración mecánica y elimina prácticamente el riesgo de coqueras en secciones delgadas y densamente armadas.

\paragraph{Metodología de reparación}
La reparación corresponde a una reconstitución del hormigón mediante mortero de reparación. El procedimiento es el siguiente:
\begin{enumerate}
    \item Demolición de la zona afectada hasta alcanzar hormigón sano, generando bordes definidos y evitando cantos en bisel que favorezcan el desprendimiento.
    \item Limpieza y preparación de la superficie, eliminando el material suelto y la lechada, y dejando una rugosidad adecuada (perfil CSP según ICRI 310.2R) que asegure la adherencia mecánica \cite{ICRI}.
    \item Aplicación de un puente de adherencia epóxico sobre la superficie limpia, dentro de su ventana de trabajabilidad y en estado \textit{tacky} al momento de recibir el mortero.
    \item Colocación de un mortero de reparación cementoso modificado con polímeros, de alta adherencia y baja permeabilidad, en los espesores por capa indicados por el fabricante.
    \item Curado del mortero según especificación, para evitar fisuración por retracción de la zona reparada.
\end{enumerate}

\subsection{Fisuración por retracción plástica en las losas}

\paragraph{Descripción y causa}
La fisuración por retracción plástica se manifiesta como fisuras superficiales, de trazado irregular, que aparecen en las primeras horas tras el vaciado y antes del fraguado inicial. Su causa es una tasa de evaporación del agua superficial superior a la tasa de exudación del hormigón, lo que genera tracciones en una masa que aún no ha desarrollado resistencia. Este problema es especialmente probable en la losa de fondo del CS y en la losa de fundación del campamento, por tratarse de grandes superficies horizontales expuestas en un clima de muy baja humedad relativa, alta radiación y vientos sostenidos. Cabe recordar que el viento suele ser más determinante que la temperatura ambiente en la aparición de estas fisuras.

\paragraph{Medida preventiva en el proceso constructivo}
La prevención pasa por controlar la evaporación superficial durante la ventana crítica. Las medidas son iniciar el curado tan pronto como sea posible tras el acabado, sin esperar al fraguado total; proteger la superficie de la incidencia directa del sol y del viento mediante toldos y barreras corta-viento; y monitorear la tasa de evaporación potencial, de modo que al superar 1,0 kg/m\textsuperscript{2}/h se implementen nebulizadores de agua y protección superficial adicional. El empleo de membrana de curado o compuesto filmógeno inmediatamente después del acabado completa la protección.

\paragraph{Metodología de reparación}
Si las fisuras son pasivas (estabilizadas, sin movimiento) y la superficie se encuentra seca, la reparación corresponde a una reparación de fisuras mediante inyección de resina epóxica, con el objetivo de devolver la monoliticidad y la transferencia de esfuerzos del elemento. El procedimiento contempla:
\begin{enumerate}
    \item Limpieza de la fisura y de su entorno, asegurando que esté seca y libre de polvo.
    \item Sellado superficial de la fisura y colocación de los puntos de inyección (\textit{packers}) cuando se utilice inyección con bomba; para fisuras accesibles desde la cara superior puede emplearse inyección por gravedad con una resina de muy baja viscosidad.
    \item Inyección de la resina epóxica, de baja viscosidad, sin retracción al curar y con buena adherencia al hormigón, hasta el llenado completo de la fisura.
    \item Retiro del sello superficial y verificación del relleno.
\end{enumerate}
En caso de fisuras de retracción muy finas y superficiales que no comprometan la sección, puede bastar con un tratamiento superficial de impermeabilización en lugar de la inyección estructural.

\subsection{Filtración a través de una junta fría en el CS}

\paragraph{Descripción y causa}
Una junta fría se produce cuando una capa de hormigón comienza a fraguar antes de recibir la siguiente, generando un plano de discontinuidad sin adherencia adecuada entre ambas. Su causa típica es una interrupción no planificada del hormigonado (falla mecánica, corte del suministro de hormigón fresco) o una mala preparación de la junta antes de reanudar el vaciado. En el muro del CS, una junta fría mal tratada constituye una vía preferente de filtración del líquido contenido, comprometiendo la función hidráulica de la estructura, en particular en la zona de mayor presión hidrostática en la base del muro.

\paragraph{Medida preventiva en el proceso constructivo}
La medida preventiva principal es planificar un hormigonado continuo, asegurando un abastecimiento ininterrumpido de hormigón fresco y aprovechando el avance del encofrado deslizante para minimizar las juntas. Si una interrupción resulta inevitable, la junta debe tratarse antes de reanudar: esperar a que el hormigón alcance resistencia suficiente, eliminar la lechada superficial mediante chorro de agua a presión o sandblasting, dejar la superficie en condición SSS y aplicar un puente de adherencia. Para la junta programada entre la losa de fondo y el muro, la incorporación de una waterstop embebida actúa como barrera física al paso del agua.

\paragraph{Metodología de reparación}
Cuando la junta presenta filtración activa, con presencia de agua, la reparación corresponde a una detención de filtraciones mediante inyección de resina de poliuretano. A diferencia de la resina epóxica, el poliuretano reacciona al contacto con el agua y expande formando una espuma que sella la vía de filtración, por lo que es apto para fisuras y juntas activas o húmedas. El procedimiento es:
\begin{enumerate}
    \item Identificación de la vía de filtración y limpieza de la junta.
    \item Perforación e instalación de \textit{packers} de inyección a lo largo de la junta.
    \item Inyección de la resina de poliuretano, que al expandir detiene el paso del agua y sella la discontinuidad.
    \item Verificación de la estanqueidad y, de ser necesario, refuerzo del tratamiento con una segunda inyección.
\end{enumerate}
Una vez detenida la filtración, si la junta requiere además recuperar continuidad estructural, puede complementarse con una reconstitución superficial mediante mortero de reparación.